% from ega0-18-fr.tex
\oldpage[0\textsubscript{IV-1}]{66}

par
\[
  x(y_1\otimes y_2\otimes\cdots\otimes y_{n+1})
  =(xy_1)\otimes y_2\otimes\cdots\otimes y_{n+1}\,;
\]
la dérivation, de degré $-1$, $d_n:\mathbf P_n\to\mathbf P_{n-1}$, est
définie par
\begin{align*}
  d_n(x_1\otimes x_2\otimes\cdots\otimes x_{n+1})
    &=(x_1x_2)\otimes x_3\otimes\cdots\otimes x_{n+1}
      -x_1\otimes(x_2x_3)\otimes\cdots\otimes x_{n+1}+\cdots\\
    &\quad+(-1)^{n-1}x_1\otimes x_2\otimes\cdots\otimes x_{n-1}
      \otimes(x_nx_{n+1})
      +(-1)^n(x_{n+1}x_1)\otimes x_2\otimes\cdots\otimes x_n,
\end{align*}
qui est bien $B$-linéaire puisque $B$ est commutatif.

Un $2$-cocycle de ce complexe à valeurs dans $L$ est une application
$B$-linéaire $h$ de $B\otimes B\otimes B$ dans $L$ telle que
$h(d_3(x\otimes y\otimes z\otimes t))=0$~; mais comme
$h(x\otimes y\otimes z)=xh(1\otimes y\otimes z)$, le cocycle $h$ est déterminé
par l'application $A$-bilinéaire
$(y,z)\longmapsto f(y,z)=h(1\otimes y\otimes z)$ de $B\times B$ dans $L$, et,
en écrivant la condition précédente pour $h$, on retrouve pour $f$ la
condition (18.4.3.1). De même, un $2$-cobord sera une application de la forme
$x\otimes y\otimes z\longmapsto h'(d_2(x\otimes y\otimes z))$, où
$h':B\otimes_A B\to L$ est linéaire, et ici encore $h'$ est déterminé par
l'application $B$-linéaire $g:y\longmapsto h'(1\otimes y)$ de $B$ dans $L$~;
on obtient alors
$h'(d_2(1\otimes x\otimes y))=xg(y)-g(xy)+yg(x)$, ce qui redonne (18.4.3.2).
On procède de même pour tout $i$, et l'on voit ainsi que l'on a
\begin{equation}
\label{0.18.4.4.1-fr}
\tag{18.4.4.1}
  \mathrm H_A^i(B,L)=\mathrm H^i(\operatorname{Hom}_B(\mathbf P_\bullet,L)).
\end{equation}
\end{env}

\begin{env}[18.4.5]
\label{0.18.4.5-fr}
Sous les conditions de (18.4.4), on peut interpréter de la même façon le
groupe $\mathrm H_A^2(B,L)^s$ (18.4.3). Pour cela, modifions le complexe
$\mathbf P_\bullet$ en degré $3$, en considérant un nouveau complexe
\[
  \mathbf P'_\bullet:\mathbf P'_3\xrightarrow{d'_3}\mathbf P_2
  \xrightarrow{d_2}\mathbf P_1\,;
\]
nous prendrons
$\mathbf P'_3=\mathbf P_3\oplus(B\otimes_A B\otimes_A B)$, $d'_3$ coïncidant
avec $d_3$ sur $\mathbf P_3$, et étant donné dans $B\otimes_A B\otimes_A B$
par
\[
  d'_3(x\otimes y\otimes z)=x\otimes y\otimes z-x\otimes z\otimes y.
\]

La relation $d_2d'_3=0$ résulte de la commutativité de $B$. Avec les notations
introduites ci-dessus, un $2$-cocycle de $\mathbf P'_\bullet$ correspond
maintenant à une application $A$-bilinéaire $f:B\times B\to L$ qui est
\emph{symétrique} et vérifie (18.4.3.1)~; par suite, on a
\[
  \mathrm H_A^2(B,L)^s
  =\mathrm H^2(\operatorname{Hom}_B(\mathbf P'_\bullet,L)).
\]
\end{env}

\begin{env}[18.4.6]
\label{0.18.4.6-fr}
Dans le cas particulier où l'on considère un \emph{corps} commutatif $k$, une
\emph{extension} $K$ de $k$, on a $\operatorname{Ext}_K^1(M,L)=0$ quels que
soient les $K$-espaces vectoriels $L$, $M$, et par suite (M, VI, 3.3 $a$)), on
a un isomorphisme canonique
\begin{equation}
\label{0.18.4.6.1-fr}
\tag{18.4.6.1}
  \mathrm H_k^i(K,L)\xrightarrow{\sim}
  \operatorname{Hom}_K(\mathrm H_i(\mathbf P_\bullet),L)
\end{equation}
et de même
\begin{equation}
\label{0.18.4.6.2-fr}
\tag{18.4.6.2}
  \mathrm H_k^2(K,L)^s\xrightarrow{\sim}
  \operatorname{Hom}_K(\mathrm H_2(\mathbf P'_\bullet),L).
\end{equation}
\end{env}

\subsection{Cas des anneaux topologiques}
\label{subsection:0.18.5-fr}

\begin{env}[18.5.1]
\label{0.18.5.1-fr}
Soient $A$, $B$ deux anneaux \emph{topologiques} dont la topologie est
\emph{linéaire}, $\rho:A\to B$ un homomorphisme continu, $L$ un $B$-bimodule
topologique, et supposons qu'il existe un idéal bilatère \emph{ouvert}
$\mathfrak R_0$ de $B$ tel que
$\mathfrak R_0.L=L.\mathfrak R_0=0$, de sorte que $L$ peut être

